% =====================================================================
%  CHAPTER 16: 健康管理与治理
% =====================================================================
\chapter{健康管理与治理}
\setcounter{chapter}{16}
\chapterintro{
掌握健康管理策略（\textbf{教学难点}）；掌握健康管理与健康治理的联系与区别；熟悉健康管理的基本步骤和核心内容；了解健康管理产生的背景、演变及现代健康管理的新特点；了解不同层次的健康管理与治理。
}

% ----- 16.1 健康管理概述 -----
\section{健康管理概述}

\subsection{健康管理产生的背景}

\begin{enumerate}[leftmargin=*]
    \item \imp{疾病谱和医学模式的转变}——慢性非传染性疾病成为主要疾病负担，要求从"治疗已病"转向"管理健康"
    \item \imp{医疗费用持续增长}——人口老龄化、新技术应用等推动医疗费用不断攀升，需要"以管理促节约"
    \item \imp{健康管理理念的兴起}——20世纪20-30年代美国最早出现健康管理萌芽（1929年Kaiser健康计划），60-70年代HMO（健康维护组织）的发展为健康管理提供了组织模式
    \item \imp{循证医学和信息技术的发展}——为健康管理的科学化和信息化提供了支撑
\end{enumerate}

\subsection{健康管理的演变}

从疾病管理→健康管理→健康治理：
\begin{itemize}[leftmargin=*]
    \item 第一阶段：\imp{疾病管理（Disease Management）}——以特定疾病患者为对象
    \item 第二阶段：\imp{健康管理（Health Management）}——扩展到健康和亚健康人群
    \item 第三阶段：\imp{健康治理（Health Governance）}——从管理到治理，强调多元主体参与合作
\end{itemize}

\subsection{健康管理的概念和内涵}

\begin{defbox}
\textbf{健康管理（Health Management）}：以健康为中心，对\keyword{个体或群体的健康危险因素}进行全面监测、分析、评估和预测，提供健康咨询和指导，并对健康危险因素进行干预的全过程。其核心是\imp{对健康风险的管理和控制}。

\textbf{核心内涵：}
\begin{enumerate}[leftmargin=*]
    \item 以健康为管理对象——而非以疾病为对象
    \item 以健康风险为核心——识别、评估和干预健康危险因素
    \item 以预防为导向——防患于未然，关注全生命周期
    \item 强调个体主动性——调动个人积极性，自我管理健康
    \item 持续动态管理——健康管理是持久、连续的过程
\end{enumerate}
\end{defbox}

\subsection{现代健康管理的新特点}

\begin{enumerate}[leftmargin=*]
    \item \imp{信息化}——电子健康档案（EHR）、移动健康（mHealth）、可穿戴设备
    \item \imp{智能化}——人工智能和大数据技术的应用
    \item \imp{个性化/精准化}——基于个体风险图谱和管理需求
    \item \imp{协同化}——多部门、多学科协同参与
    \item \imp{全生命周期化}——从出生到死亡的全程健康管理
\end{enumerate}

% ----- 16.2 健康管理基本内容 -----
\section{健康管理的基本内容与步骤}

\subsection{健康管理的基本步骤}

\begin{keybox}
\textbf{健康管理的四个基本步骤：}

\begin{enumerate}[leftmargin=*]
    \item \imp{第一步：健康信息收集与监测}
    \begin{itemize}
        \item 个人基本信息（年龄、性别、家族史等）
        \item 健康体检信息（体格检查、实验室检查等）
        \item 行为生活方式信息（吸烟、饮酒、饮食、运动等）
        \item 心理与社会状况评估
    \end{itemize}

    \item \imp{第二步：健康风险评估（HRA）}
    \begin{itemize}
        \item 评估个体或群体未来发生特定疾病或健康不良后果的风险
        \item 识别主要可改变的健康危险因素
        \item 确定干预的优先顺序
    \end{itemize}

    \item \imp{第三步：健康管理计划制定}
    \begin{itemize}
        \item 设定健康改善目标（可测量、可达到、有时间限制）
        \item 制定个性化的干预方案（行为改善、治疗和管理方案）
        \item 明确实施方式和时间安排
    \end{itemize}

    \item \imp{第四步：健康管理干预与评估}
    \begin{itemize}
        \item 实施健康干预（健康教育、行为矫正、临床干预等）
        \item 定期随访和监测
        \item 效果评估和计划调整
    \end{itemize}
\end{enumerate}
\end{keybox}

\subsection{健康管理策略（教学难点*）}

\begin{keybox}
\imp{（教学难点*）}\textbf{健康管理的六大基本策略：}

\begin{enumerate}[leftmargin=*]
    \item \imp{生活方式管理（Lifestyle Management）}——通过健康教育和行为干预，帮助人们采纳和维持健康的生活方式。核心手段：健康教育、激励机制、行为矫正。目标：减少健康危险行为（吸烟、酗酒、不良饮食、缺乏运动等）

    \item \imp{需求管理（Demand Management）}——运用远程医疗、健康咨询、自我管理支持等手段，帮助人们更好地利用卫生服务，避免不必要和不恰当的卫生服务使用。主要工具：24小时电话健康咨询、在线健康信息、分诊系统

    \item \imp{疾病管理（Disease Management）}——以特定疾病患者为对象，采用系统化、协调化的服务模式，改善疾病结局。强调：循证医学指导、患者自我管理、多学科团队协作、持续质量改进。典型对象：糖尿病、高血压、哮喘、COPD等慢性病

    \item \imp{灾难性伤病管理（Catastrophic Case Management）}——针对严重伤病或高额医疗费用病例的专门管理，通过个案管理和专业评估，在确保医疗质量的同时控制费用。要求：专业的个案管理师、综合评估和资源配置

    \item \imp{残疾管理（Disability Management）}——减少工作场所残疾发生，促进因伤病缺勤者尽早康复和重返工作。关注：早期康复介入、工作环境调整、心理支持和职业康复

    \item \imp{综合人群健康管理（Population Health Management）}——整合以上多种策略，对整个人群的健康进行系统化综合管理，强调：数据驱动的风险评估和分层管理、多策略的协调实施、跨部门协作、健康结果和成本效益评估
\end{enumerate}
\end{keybox}

% ----- 16.3 健康治理 -----
\section{健康治理}

\begin{defbox}
\textbf{健康治理（Health Governance）}：政府、市场、社会组织、公民等\keyword{多元主体}通过协商、合作和共同行动，共同应对健康问题和挑战的持续过程。

\textbf{健康治理的特点：}
\begin{enumerate}[leftmargin=*]
    \item \imp{多元主体参与}——不仅限于政府，还包括企业、社会组织、社区、公民等
    \item \imp{协商共治}——各主体间通过协商而非命令来达成共识
    \item \imp{透明和问责}——决策过程透明、各主体有明确责任并接受监督
    \item \imp{网络化治理}——治理结构从"垂直层级"转向"网络协作"
    \item \imp{全球视野}——健康问题超越国界，需要全球治理机制
\end{enumerate}
\end{defbox}

\subsection{健康管理与健康治理的联系与区别}

\begin{keybox}
\textbf{联系：}
\begin{itemize}[leftmargin=*]
    \item 健康治理是健康管理理念的延伸和升级
    \item 两者都强调以健康为中心、预防为主
    \item 健康治理为健康管理提供了更广泛的制度框架和组织支撑
\end{itemize}

\textbf{区别：}
\centering
\textbf{表16-1 健康管理与健康治理的比较}
\bigskip

\begin{tabularx}{\textwidth}{lXX}
\hline
\textbf{比较维度} & \textbf{健康管理} & \textbf{健康治理} \\
\hline
主体 & 以卫生专业人员和管理者为主 & 多元主体（政府、市场、社会、公民） \\
方式 & 以专业手段进行管理和干预 & 以协商、合作、共同行动实现 \\
对象 & 个体或群体的健康风险 & 健康的社会决定因素和系统问题 \\
范围 & 侧重卫生服务领域 & 跨部门、跨领域的社会性治理 \\
机制 & 计划、组织、协调、控制 & 参与、协商、透明、问责 \\
\hline
\end{tabularx}
\bigskip
\end{keybox}

% ----- 16.4 不同层次的健康管理与治理 -----
\section{不同层次的健康管理与治理}

\begin{enumerate}[leftmargin=*]
    \item \imp{以社区为基础的健康管理与治理}——社区是健康管理和治理的基础单元。家庭医生签约服务、社区网格化管理、健康社区建设
    \item \imp{卫生系统健康管理与治理}——医保支付方式改革（按人头付费、DRG等）、分级诊疗、医联体、信息化建设
    \item \imp{国家健康管理与治理}——制定国家健康战略（如"健康中国2030"）、健康立法、多部门协调机制（爱国卫生运动委员会）
    \item \imp{全球健康治理}——WHO在全球健康治理中的核心作用、国际卫生条例（IHR）、全球基金（GFATM）、新发传染病的全球应对机制
\end{enumerate}

% ----- 复习题 -----
\section{复习题}

\subsection*{一、选择题}
\begin{enumerate}[leftmargin=*, itemsep=3pt]
    \item 健康管理的核心是：
    \begin{enumerate}
        \item[(A)] 疾病治疗
        \item[(B)] 健康风险的管理和控制
        \item[(C)] 药品管理
        \item[(D)] 医院管理
    \end{enumerate}
    \item 以下哪项不是健康管理的四大基本步骤之一？
    \begin{enumerate}
        \item[(A)] 健康信息收集与监测
        \item[(B)] 健康风险评估
        \item[(C)] 疾病诊断与治疗
        \item[(D)] 健康管理干预与评估
    \end{enumerate}
    \item 健康管理与健康治理的最主要区别在于：
    \begin{enumerate}
        \item[(A)] 前者是治疗，后者是预防
        \item[(B)] 前者以专业人员为主，后者强调多元主体参与
        \item[(C)] 前者关注个体，后者不关注个体
        \item[(D)] 两者没有区别
    \end{enumerate}
    \item 健康管理中"灾难性伤病管理"的主要对象是：
    \begin{enumerate}
        \item[(A)] 所有慢性病患者
        \item[(B)] 健康人群
        \item[(C)] 严重伤病或高额医疗费用病例
        \item[(D)] 亚健康人群
    \end{enumerate}
\end{enumerate}

\subsection*{二、名词解释}
\begin{enumerate}[leftmargin=*, itemsep=3pt]
    \item \textbf{健康管理（Health Management）}
    \item \textbf{健康治理（Health Governance）}
    \item \textbf{生活方式管理}
\end{enumerate}

\subsection*{三、简答题}
\begin{enumerate}[leftmargin=*, itemsep=3pt]
    \item 简述健康管理的六大基本策略。
    \item 简述健康管理的基本步骤。
    \item 试述健康管理与健康治理的联系与区别。
    \item 简述现代健康管理的新特点。
\end{enumerate}

\subsection*{参考答案要点}

\subsubsection*{一、选择题}
1. (B)；2. (C)；3. (B)；4. (C)

\subsubsection*{二、名词解释}
\begin{enumerate}[leftmargin=*]
    \item \textbf{健康管理}：对个体或群体的健康危险因素进行全面监测、分析、评估和干预的全过程，核心是对健康风险的管理和控制。
    \item \textbf{健康治理}：多元主体（政府、市场、社会、公民）通过协商合作共同应对健康问题和挑战的持续过程。
    \item \textbf{生活方式管理}：通过健康教育和行为干预帮助人们采纳和维持健康生活方式的策略，核心手段包括健康教育、激励机制和行为矫正。
\end{enumerate}

\subsubsection*{三、简答题}
\begin{enumerate}[leftmargin=*]
    \item 六大策略：①生活方式管理（健康教育+行为干预）；②需求管理（减少不必要服务使用）；③疾病管理（慢性病系统化管理）；④灾难性伤病管理（严重病例个案管理）；⑤残疾管理（促进康复和重返工作）；⑥综合人群健康管理（数据驱动、多策略协调）。
    \item 四步：信息收集监测→健康风险评估（HRA）→制定管理计划→干预实施与效果评估。
    \item 联系：治理是管理的延伸升级，都以健康为中心。区别：管理以专业人员为主体，侧重技术手段和风险管控；治理以多元主体参与为特征，侧重协商合作和制度建设，范围更广泛（跨部门vs.卫生服务领域）。
    \item 五大新特点：信息化（EHR、mHealth）、智能化（AI、大数据）、个性化/精准化、协同化（多部门多学科）、全生命周期化。
\end{enumerate}

\newpage
